\documentclass[11pt,notitlepage]{article}

\usepackage{pslatex}
\usepackage[T1]{fontenc}
\usepackage[english]{babel}
\usepackage{vmargin}
\usepackage[utf8]{inputenc}
\usepackage{setspace}
\usepackage{url}
\usepackage{graphicx,float}
\usepackage{amssymb, amsmath}
\usepackage{chngcntr}
\usepackage{hyperref}
\thispagestyle{empty}
\floatstyle{plain}\restylefloat{figure}

\hypersetup{
	pdftitle={Monetary Policy Review: Press Remarks},
	pdfauthor={Martien Lubberink},
	colorlinks=true,
	linkcolor=blue,
	citecolor=blue,
	filecolor=blue,
	urlcolor=blue
}

\onehalfspacing

\setmarginsrb{1in}{1in}{1in}{.5in}{11pt}{11pt}{11pt}{33pt}
\newcommand{\clearemptydoublepage}{\newpage{\pagestyle{empty}\cleardoublepage}}

\makeatletter
\def\url@leostyle{%
	\@ifundefined{selectfont}{\def\UrlFont{\sf}}{\def\UrlFont{\itshape}}}
\makeatother
\urlstyle{leo}

\newcounter{parnum}
\newcommand{\N}{%
	\noindent\refstepcounter{parnum}%
	\makebox[\parindent][l]{\textbf{\arabic{parnum}.}}}
\setlength{\parindent}{1.6em}

% --- pandoc support macros ---
\providecommand{\tightlist}{%
  \setlength{\itemsep}{0pt}\setlength{\parskip}{0pt}}
\providecommand{\passthrough}[1]{#1}

\begin{document}
\textbf{The review finds that the monetary policy framework had become more flexible before COVID, and that this increased its vulnerability to policy mistakes under extreme uncertainty.} The reviewers conclude that the evolution of inflation targeting, culminating in the 2018 Dual Mandate, added scope for discretion and placed less emphasis on price stability. They do not say that the initial COVID response was wrong: they describe the early monetary response as rapid, innovative and reasonable given the information available at the time. Their concern is that the framework was not sufficiently robust once the economic outlook changed.

\textbf{The central problem identified by the review was not the decision to provide strong support in the initial shock, but the failure to adjust that support promptly as the recovery proved stronger than expected.} The review finds that forecasts were persistently too pessimistic, incoming positive information was discounted, and policy remained highly accommodative for too long. It concludes that the resulting overheating contributed to inflation reaching 7.3 percent and that the subsequent tightening imposed substantial costs on growth and employment.

\textbf{The review explicitly identifies the Dual Mandate as one of the institutional changes that weakened the guardrails around monetary policy.} Its conclusion is that the 2018 legislative change added scope for discretion and guided the Monetary Policy Committee to focus less on price stability. The review therefore recommends strengthening the framework and, specifically, says that policy should be more systematic and rely less on discretion. That is a finding about institutional design, not a judgement about the motives or abilities of the people operating the framework.

\textbf{The review also makes clear that this was not simply a story about forecasting errors or the extraordinary nature of COVID.} Forecast errors are unavoidable in a crisis, and the reviewers acknowledge that other central banks made mistakes. But they find that New Zealand's errors were unusually large and persistently one-sided, and argue that vulnerabilities in the policy strategy were a major factor. Their recommendations include greater use of scenarios, simple policy rules, greater attention to near-term inflation signals and more open consideration of alternative policy paths.

\textbf{The lesson for future governments is that monetary policy needs strong institutional guardrails precisely when the outlook is most uncertain.} The review says that price stability remains the most valuable contribution monetary policy can make to economic progress and welfare, while recognising the importance of sustainable employment. My focus is therefore on what the review recommends for the future: a more robust and systematic framework, clearer accountability, better stress-testing of policy choices and arrangements that reduce the risk that a reasonable response to an extraordinary shock becomes excessive and difficult to unwind.

\end{document}
